\chapter*{Glosario de términos}
\addcontentsline{toc}{chapter}{Glosario de términos}
\markboth{GLOSARIO DE TÉRMINOS}{}

\begin{description}
\item[OWASP] Open Worldwide Application Security Project. Fundación sin
  ánimo de lucro que produce, de forma abierta y colaborativa, herramientas,
  guías y catálogos de referencia en seguridad de aplicaciones, entre ellos
  el OWASP Top~10 y la Web Security Testing Guide (WSTG).

\item[WSTG] Web Security Testing Guide. Guía de la OWASP Foundation que
  descompone la auditoría de una aplicación web en fases (gestión de
  identidad, autenticación, autorización, gestión de sesión, validación de
  entradas, etc.) y numera una prueba concreta para cada tipo de fallo,
  usada en este trabajo para situar cada vulnerabilidad de \breachlab
  dentro de un marco de referencia reconocido en la industria.

\item[PTES] Penetration Testing Execution Standard. Marco de referencia para
  estructurar una auditoría de seguridad ofensiva de alcance general (no
  específico de aplicaciones web), citado en este trabajo como alternativa
  considerada frente a la OWASP WSTG.

\item[OSSTMM] Open Source Security Testing Methodology Manual. Manual de
  metodología de pruebas de seguridad de infraestructura, orientado
  igualmente a un alcance más general que el de la OWASP WSTG.

\item[NIST SP 800-115] Guía técnica del National Institute of Standards and
  Technology para la planificación y ejecución de pruebas y evaluaciones de
  seguridad de la información.

\item[MITRE ATT\&CK] Catálogo público, mantenido por MITRE, de tácticas y
  técnicas empleadas por atacantes reales, organizado por fases de un
  ataque; se menciona en este trabajo como marco de referencia para el
  encadenamiento de vulnerabilidades del caso de estudio.

\item[CWE] Common Weakness Enumeration. Catálogo mantenido por MITRE que
  clasifica y describe tipos genéricos de debilidades de software (por
  ejemplo, CWE-89 para la inyección SQL), independientemente de una
  vulnerabilidad concreta.

\item[CVSS] Common Vulnerability Scoring System. Estándar mantenido por
  FIRST.org para expresar de forma cuantitativa la severidad de una
  vulnerabilidad mediante un vector de métricas y una puntuación de 0 a 10.

\needspace{15\baselineskip}
\item[SQL Injection (inyección SQL)] Vulnerabilidad que permite alterar la
  lógica de una consulta a base de datos insertando código SQL a través de
  una entrada no validada ni parametrizada.

\item[XSS (Cross-Site Scripting)] Vulnerabilidad que permite ejecutar
  script arbitrario en el navegador de otro usuario. Puede ser reflejado,
  almacenado o basado en DOM, según dónde se materialice el fallo.

\item[XSS basado en DOM] Variante de XSS en la que el payload nunca pasa por
  el servidor: se inserta directamente en el árbol DOM desde JavaScript
  ejecutado en el navegador, típica de aplicaciones SPA.

\item[CSRF (Cross-Site Request Forgery)] Vulnerabilidad que aprovecha que el
  navegador adjunta automáticamente las cookies de sesión para forzar a la
  víctima a enviar, sin saberlo, una petición a una aplicación en la que
  está autenticada.

\item[IDOR (Insecure Direct Object Reference)] Categoría de control de
  acceso roto en la que un identificador de objeto (un id de nota, de
  pedido, etc.) puede manipularse para acceder a recursos de otro usuario
  sin autorización.

\item[SPA (Single Page Application)] Patrón de arquitectura web en el que
  el navegador carga una única página que actualiza su contenido de forma
  dinámica, habitualmente consumiendo una API independiente.

\item[CORS (Cross-Origin Resource Sharing)] Mecanismo del navegador que
  restringe las peticiones entre orígenes distintos; una
  aplicación que sirve frontend y backend en orígenes diferentes debe
  habilitarlo explícitamente, lo que introduce restricciones adicionales
  que \breachlab evita sirviendo ambos componentes bajo el mismo origen.

\item[Cookie HttpOnly] Atributo de una cookie que impide su lectura desde
  JavaScript, reduciendo el impacto de un XSS al no poder exfiltrarse
  mediante script.

\item[Cookie SameSite] Atributo de una cookie que restringe en qué contexto
  de navegación de terceros se envía. El valor \texttt{Lax} la envía en
  navegaciones de nivel superior por GET, pero no en subrecursos
  cross-site.

\item[CSP (Content-Security-Policy)] Cabecera HTTP que permite restringir
  qué orígenes pueden servir script, estilos u otros recursos en una
  página, mitigando el impacto de un XSS aunque el saneado del contenido
  falle.

\item[Hash de contraseña] Resultado de aplicar una función criptográfica
  de un solo sentido (en este trabajo, mediante Werkzeug) a una contraseña,
  de forma que pueda verificarse sin almacenarla en texto plano.

\needspace{15\baselineskip}
\item[Fijación de sesión (\emph{session fixation})] Ataque en el que un
  identificador de sesión obtenido antes de que la víctima se autentique se
  reutiliza después con sus privilegios; se evita invalidando la sesión
  (por ejemplo, con \texttt{session.clear()}) en el momento del inicio de
  sesión.

\item[Pentesting (test de penetración)] Ejercicio de seguridad ofensiva
  autorizado que simula un ataque real contra un sistema con el fin de
  identificar y demostrar vulnerabilidades explotables.

\item[DAST (Dynamic Application Security Testing)] Técnica de análisis de
  seguridad que examina una aplicación en ejecución, sin acceso a su código
  fuente, enviándole peticiones y observando sus respuestas.

\item[Clickjacking] Técnica que consiste en embeber una página legítima en
  un iframe invisible o disfrazado para engañar al usuario y hacer que
  interactúe con ella sin darse cuenta.
\end{description}
